\chapter{Đặt vấn đề và tổng quan}

\section{Đặt vấn đề}
Nông nghiệp từ lâu giữ vai trò quan trọng trong quá trình phát triển kinh tế - xã hội, bảo đảm an ninh lương thực và tạo sinh kế cho người dân Việt Nam \cite{world-bank-vietnam-future-jobs}. 
Những số liệu gần đây cho thấy quy mô và năng lực sản xuất của nông nghiệp Việt Nam vẫn ở mức đáng kể. 
Năm 2024, tổng sản lượng lương thực đạt khoảng 48,1 triệu tấn, trong đó sản lượng lúa đạt 43,7 triệu tấn, với năng suất bình quân khoảng 61,4 tạ/ha \cite{ppd-crop-production-79-years}. 
Cũng trong năm 2024, kim ngạch xuất khẩu nông, lâm, thủy sản đạt mức kỷ lục 62,5 tỷ USD, trong đó riêng các sản phẩm trồng trọt đóng góp khoảng 32,8 tỷ USD và xuất khẩu nông sản khoảng gần 20 tỷ USD \cite{ppd-crop-production-79-years}. 
Mặc dù ngành nông nghiệp hiện chỉ chiếm khoảng 10\% GDP, hoạt động của ngành vẫn gắn với khoảng 60\% dân cư sinh sống tại khu vực nông thôn, cho thấy vai trò quan trọng đối với sinh kế, an ninh lương thực và hoạt động xuất khẩu của Việt Nam \cite{mae-agriculture-2025}.

Trong tương lai, nhu cầu nâng cao năng suất nông nghiệp được dự báo tiếp tục gia tăng \cite{fao-land-water-2025}.
Bên cạnh đó, FAO cũng cảnh báo  khoảng 60\% diện tích đất sẽ bị suy thoái do tác động của con người trên đất nông nghiệp, vì vậy mở rộng diện tích canh tác không còn là giải pháp bền vững.\cite{fao-land-water-2025}
Trong bối cảnh tài nguyên đất và nước chịu nhiều áp lực, việc nâng cao năng suất trên diện tích hiện có, sử dụng tài nguyên hiệu quả và quản lý dinh dưỡng hợp lý sẽ trở thành những yêu cầu ngày càng quan trọng đối với nền nông nghiệp tương lai.
Cùng với quan ngại này, định hướng nông nghiệp Việt Nam không chỉ tập trung vào gia tăng sản lượng mà còn hướng tới nâng cao chất lượng, hiệu quả kinh tế, tiết kiệm vật tư đầu vào và giảm tác động đến môi trường \cite{mae-agriculture-2025}. 

Trong bối cảnh đó, ngành sản xuất phân bón giữ vị trí quan trọng trong chuỗi cung ứng đầu vào của sản xuất nông nghiệp. 
Theo thông tin của Bộ Nông nghiệp và Môi trường năm 2025, Việt Nam hiện sản xuất khoảng 20,9 triệu tấn phân bón mỗi năm, trong đó khoảng 3 triệu tấn là phân bón hữu cơ \cite{mae-fertilizer-production}. 
Cơ cấu sản phẩm phân bón trong nước cũng đang có xu hướng chuyển dịch theo hướng đa dạng hóa và ưu tiên các sản phẩm thân thiện với môi trường; tỷ lệ sản phẩm phân bón hữu cơ được phép lưu hành hiện chiếm khoảng 27\% tổng số sản phẩm phân bón được phép lưu hành \cite{mae-one-million-hectares}. 
Đồng thời, định hướng quản lý hiện nay khuyến khích sử dụng phân bón tiết kiệm, cân đối và hiệu quả, tăng sử dụng phân bón hữu cơ và từng bước giảm sự phụ thuộc vào phân bón vô cơ \cite{mae-one-million-hectares}.

Trong quá trình nâng cao năng suất cây trồng, phân bón là một trong những yếu tố đầu vào có vai trò đặc biệt quan trọng. 
Phân bón bổ sung các nguyên tố dinh dưỡng mà đất không thể cung cấp đầy đủ cho cây trồng, qua đó hỗ trợ duy trì năng suất và chất lượng nông sản.
IFA hiện ước tính khoảng một nửa lượng thực phẩm được sản xuất trên thế giới có sự đóng góp của phân bón khoáng, cho thấy ngành sản xuất phân bón có mối liên hệ trực tiếp với khả năng duy trì sản lượng nông nghiệp và an ninh lương thực.\cite{ifa-food-nutrition-security}
Cùng với sự gia tăng nhu cầu sản xuất nông nghiệp, ngành phân bón thế giới tiếp tục duy trì quy mô thị trường lớn và được dự báo còn tăng trưởng trong những năm tới.\cite{grand-view-fertilizer-market}
Cơ cấu của ngành phân bón cũng đang có xu hướng chuyển dịch từ việc đơn thuần tăng lượng phân sử dụng sang tăng hiệu quả sử dụng dinh dưỡng, đa dạng hóa công thức và phát triển các sản phẩm chuyên biệt.\cite{fortune-organic-fertilizer-market}
Các xu hướng này cho thấy thị trường đang ngày càng quan tâm đến những sản phẩm không chỉ cung cấp N, P, K mà còn bổ sung trung lượng, vi lượng, thành phần hữu cơ và các giải pháp giúp nâng cao hiệu quả sử dụng dinh dưỡng.\cite{fortune-organic-fertilizer-market}

Đối với sự đa dạng ngày càng lớn của các sản phẩm phân bón, việc kiểm soát chất lượng trở nên vô cùng quan trọng, đặc biệt về thành phần và hàm lượng dinh dưỡng của sản phẩm. 
Những sai lệch trong nguyên liệu, tỷ lệ phối trộn hoặc quá trình sản xuất có thể làm chất lượng thành phẩm không đạt yêu cầu, ảnh hưởng trực tiếp đến hiệu quả sử dụng trên cây trồng, dư lượng hoá chất trên thực phẩm, sức khoẻ người tiêu dùng và môi trường. 
Vì vậy, việc kiểm soát cần được thực hiện xuyên suốt từ nguyên liệu đầu vào, bán thành phẩm đến thành phẩm cuối cùng, trong đó phòng thí nghiệm kiểm định chất lượng giữ vai trò cung cấp kết quả phân tích làm cơ sở đánh giá và quyết định chất lượng sản phẩm.

Vì vậy, công tác kiểm nghiệm cần được thực hiện xuyên suốt từ nguyên liệu đầu vào, bán thành phẩm trong quá trình sản xuất đến thành phẩm cuối dây chuyền. 
Trong đó, phòng thí nghiệm kiểm định chất lượng giữ vai trò thực hiện lấy mẫu, xử lý và phân tích mẫu, cung cấp kết quả kiểm nghiệm làm cơ sở đánh giá chất lượng và phát hiện những sai lệch có thể phát sinh trong quá trình sản xuất. 
Xuất phát từ thực tế này, quá trình thực tập tại phòng thí nghiệm kiểm định chất lượng sản phẩm của Công ty TNHH Quốc tế Khánh Sinh là cơ hội để tìm hiểu trực tiếp quy trình lấy mẫu, xử lý mẫu, thực hiện các phép phân tích và kiểm soát chất lượng phân bón trước khi sản phẩm được đưa ra tiêu thụ trên thị trường.

\section{Mục đích môn học}
Các mục tiêu cụ thể của kỳ thực tập được xác định như sau:

Củng cố và vận dụng các kiến thức đã học về hóa phân tích, kiểm nghiệm và kiểm soát chất lượng vào công việc thực tế tại phòng thí nghiệm của doanh nghiệp sản xuất phân bón.

Học hỏi và tích lũy kinh nghiệm trong môi trường phòng thí nghiệm thực tế, rèn luyện các kỹ năng như lấy mẫu, chuẩn bị mẫu, pha hóa chất, sử dụng thiết bị phân tích và xử lý các tình huống phát sinh trong quá trình kiểm nghiệm.

Rèn luyện tác phong làm việc cẩn thận, chính xác và tuân thủ quy trình; đồng thời nâng cao kỹ năng phối hợp, giao tiếp và làm việc với các bộ phận liên quan như Phòng Kỹ thuật và Bộ phận Sản xuất.

Tìm hiểu và tham gia vào quy trình kiểm soát chất lượng phân bón từ nguyên liệu đầu vào, bán thành phẩm đến thành phẩm cuối cùng; thu thập số liệu thực tế phục vụ cho báo cáo thực tập, luận văn tốt nghiệp và định hướng công việc sau khi ra trường.

Làm quen với hệ thống thiết bị và phương pháp kiểm nghiệm đang được sử dụng tại doanh nghiệp, qua đó nâng cao khả năng ứng dụng kiến thức chuyên ngành vào hoạt động kiểm soát chất lượng trong sản xuất phân bón.

\section{Tổng quan về phân bón}
\subsection{Khái quát về phân bón và dinh dưỡng cây trồng}
Phân bón là một trong những nguồn bổ sung dinh dưỡng quan trọng cho cây trồng, được sử dụng nhằm cung cấp những nguyên tố mà đất không thể đáp ứng đầy đủ theo nhu cầu sinh trưởng và năng suất của cây. 
Ba nguyên tố dinh dưỡng khoáng chính đối với cây trồng là nitơ (N), phospho (P) và kali (K), đây cũng là ba nguyên tố được sử dụng phổ biến nhất trong sản xuất phân bón thương mại. \cite{fao-plant-nutrition,fao-fertilizer-specifications}.

Bên cạnh N, P và K, cây trồng còn cần các nguyên tố đa lượng thứ cấp gồm canxi (Ca), magie (Mg) và lưu huỳnh (S), cùng các nguyên tố vi lượng như bo (B), sắt (Fe), mangan (Mn), đồng (Cu), kẽm (Zn) và molypden (Mo) \cite{fao-plant-nutrition}. 
Việc phân biệt đa lượng và vi lượng dựa chủ yếu vào lượng cây cần hấp thu, không phản ánh mức độ quan trọng về sinh lý của từng nguyên tố. \cite{fao-plant-nutrition}
Theo FAO, N và K chiếm khoảng 80\% tổng lượng dinh dưỡng khoáng trong cây, P, S, Ca và Mg chiếm khoảng 19\%, trong khi tổng các nguyên tố vi lượng chiếm dưới 1\% \cite{fao-plant-nutrition}.
Tuy nhiên, sinh trưởng và năng suất của cây phụ thuộc vào sự cân bằng giữa các chất dinh dưỡng, bởi sự thiếu hụt một nguyên tố thiết yếu có thể trở thành yếu tố giới hạn ngay cả khi các nguyên tố khác được cung cấp đầy đủ.  
Đây chính là cơ sở của nguyên lý dinh dưỡng cân đối và cũng giải thích vì sao các công thức phân bón hiện đại ngày càng được thiết kế với nhiều thành phần dinh dưỡng khác nhau.

\subsection{Phân đơn}
Phân đơn (\textit{straight fertilizer}) là nhóm phân bón chủ yếu cung cấp một trong ba nguyên tố dinh dưỡng chính N, P hoặc K theo cách phân loại của FAO \cite{fao-fertilizer-specifications}.
Phân đơn có thể được sử dụng trực tiếp để bổ sung một nguyên tố cụ thể hoặc làm nguyên liệu phối trộn thành các công thức NP, NK và NPK \cite{fao-fertilizer-specifications}.

\subsubsection{Phân đạm - phân N}
Phân đạm có chức năng chính là cung cấp nitơ cho cây trồng. 
Nitơ là thành phần của diệp lục và protein, đồng thời tham gia mạnh vào quá trình hình thành thân, lá, chồi và sinh khối. 
Cây chủ yếu hấp thu nitơ từ đất dưới dạng nitrat (\(\mathrm{NO_3^-}\)) và amoni (\(\mathrm{NH_4^+}\)).
Khi cây thiếu nitơ, tốc độ sinh trưởng giảm, diện tích lá nhỏ, khả năng đẻ nhánh giảm và lá thường xuất hiện hiện tượng vàng do hàm lượng chlorophyll giảm. 
Ngược lại, việc cung cấp nitơ quá mức có thể kéo dài giai đoạn sinh trưởng sinh dưỡng và làm chậm quá trình trưởng thành của cây. \cite{fao-plant-nutrition}

Một trong những loại phân đạm phổ biến nhất là urea, \(\mathrm{CO(NH_2)_2}\).
Theo thông số kỹ thuật của FAO, urê thương mại có hàm lượng nitơ tổng số tối thiểu khoảng 46\% theo khối lượng chất khô, do đó đây là nguồn phân đạm có hàm lượng dinh dưỡng cao \cite{fao-fertilizer-specifications}.
Một nguồn phân đạm khác là ammonium sulfate, trong đó FAO quy định hàm lượng N tối thiểu khoảng 20\% và đồng thời chứa tối thiểu khoảng 23\% S. \cite{fao-fertilizer-specifications}
Vì vậy, một số nguồn phân đạm không chỉ cung cấp N mà còn có thể đóng vai trò cung cấp thêm các nguyên tố dinh dưỡng thứ cấp cho cây trồng.

Đối với hoạt động kiểm định chất lượng, hàm lượng nitơ tổng số là một trong những chỉ tiêu quan trọng nhất của nhóm phân đạm, bởi sự sai lệch hàm lượng N sẽ làm thay đổi trực tiếp lượng dinh dưỡng thực tế được cung cấp cho cây.

\subsubsection{Phân lân - phân P}
Phân lân là nhóm phân bón được sử dụng chủ yếu để bổ sung phospho (P) cho cây trồng. 
Trong hệ thống công bố hàm lượng phân bón, phospho thường không được biểu thị trực tiếp dưới dạng P mà được quy đổi và ghi dưới dạng \(\mathrm{P_2O_5}\).\cite{fao-fertilizer-specifications}

Cây hấp thu phospho chủ yếu dưới các dạng \(\mathrm{H_2PO_4^-}\) và \(\mathrm{HPO_4^{2-}}\), trong đó tỷ lệ giữa hai dạng phụ thuộc vào pH của môi trường đất.
Phospho tham gia vào quá trình phân chia tế bào, kéo dài rễ, phát triển hạt và quả, đồng thời các hợp chất chứa P như ATP và ADP đóng vai trò quan trọng trong quá trình trao đổi và vận chuyển năng lượng trong tế bào.
Thiếu phospho có thể làm hạn chế sinh trưởng, giảm phát triển của bộ rễ, giảm khả năng đẻ nhánh và làm chậm quá trình trưởng thành của cây. 
So với khả năng di chuyển tương đối tốt bên trong cây, phospho có tính di động thấp hơn trong đất, vì vậy khả năng cung cấp P cho bộ rễ chịu ảnh hưởng đáng kể bởi tính chất đất và điều kiện môi trường.

Các loại phân lân đơn phổ biến gồm supe lân đơn (Single Superphosphate - SSP) và supe lân ba (Triple Superphosphate - TSP). 
Theo thông số kỹ thuật tham khảo của FAO, SSP có hàm lượng \(\mathrm{P_2O_5}\) hòa tan trong nước tối thiểu khoảng 15,8\% và đồng thời cung cấp khoảng 11\% S, trong khi TSP có hàm lượng \(\mathrm{P_2O_5}\) tổng số tối thiểu khoảng 46\%, trong đó \(\mathrm{P_2O_5}\) hòa tan trong nước tối thiểu khoảng 42,5\% \cite{fao-fertilizer-specifications}.

Vì vậy, đối với phân lân, việc chỉ xác định tổng hàm lượng phospho chưa phải lúc nào cũng phản ánh đầy đủ đặc tính của sản phẩm. 
Tuỳ từng loại phân, chúng ta sẽ có các yêu cầu chỉ tiêu phân tích khác nhau nhằm đánh giá chất lượng sản phẩm. \cite{fao-fertilizer-specifications}

\subsubsection{Phân kali - phân K}
Phân kali là nhóm phân bón được sử dụng để cung cấp kali (K) cho cây trồng, và hàm lượng kali trong sản phẩm thương mại thường được biểu thị quy đổi dưới dạng \(\mathrm{K_2O}\).\cite{fao-fertilizer-specifications}

Cây hấp thu kali dưới dạng ion \(\mathrm{K^+}\), và đây là nguyên tố khoáng có hàm lượng lớn thứ hai trong cây sau nitơ.
Kali tham gia vào hoạt động của hơn 60 hệ enzyme, quá trình quang hợp, vận chuyển sản phẩm quang hợp đến các cơ quan dự trữ và điều hòa trạng thái nước của cây. 
Kali cũng có liên quan đến khả năng chống chịu của cây trước một số điều kiện bất lợi như khô hạn, lạnh, sâu bệnh và hiện tượng đổ ngã.

Hai nguồn phân kali phổ biến là potassium chloride (\(\mathrm{KCl}\), MOP) và potassium sulfate (\(\mathrm{K_2SO_4}\), SOP).
Theo thông số kỹ thuật của FAO, potassium chloride thương mại có hàm lượng kali hòa tan quy đổi theo \(\mathrm{K_2O}\) tối thiểu khoảng 60\%, trong khi potassium sulfate có hàm lượng \(\mathrm{K_2O}\) tối thiểu khoảng 50\% và đồng thời chứa tối thiểu khoảng 17,5\% S.\cite{fao-fertilizer-specifications}

Trong sản xuất phân bón phối trộn, \(\mathrm{KCl}\) là một trong những nguyên liệu có thể được sử dụng để điều chỉnh hàm lượng kali của các công thức NPK và NK theo yêu cầu của từng sản phẩm.
Đối với phân kali, việc xác định hàm lượng \(\mathrm{K_2O}\) là chỉ tiêu quan trọng để đánh giá chất lượng sản phẩm. \cite{fao-fertilizer-specifications}

\subsection{Phân đa dinh dưỡng}
Trong phạm vi báo cáo này, phân đa dinh dưỡng được hiểu là các sản phẩm cung cấp từ hai nguyên tố dinh dưỡng đa lượng chính trở lên, chẳng hạn NP, NK hoặc NPK, đồng thời có thể được bổ sung thêm các nguyên tố đa lượng thứ cấp và vi lượng.

Theo FAO, phân compound hoặc complex chứa ít nhất hai trong ba nguyên tố dinh dưỡng chính N, P và K, khác với phân đơn chỉ cung cấp một trong ba nguyên tố này. \cite{fao-fertilizer-specifications}
Các công thức NPK được biểu thị theo thứ tự N - \(\mathrm{P_2O_5}\) - \(\mathrm{K_2O}\), trong đó ba con số trên nhãn thể hiện tỷ lệ phần trăm khối lượng của các thành phần dinh dưỡng tương ứng. \cite{fao-fertilizer-specifications}
Việc thay đổi tỷ lệ N-P-K cho phép xây dựng các công thức phân bón phù hợp với nhu cầu dinh dưỡng khác nhau của từng loại cây và từng giai đoạn sinh trưởng. 

Ngoài ba nguyên tố N, P và K, các công thức phân đa dinh dưỡng có thể được bổ sung Ca, Mg và S, là những nguyên tố đa lượng thứ cấp cần thiết đối với cây trồng. \cite{fao-integrated-nutrition}
Canxi tham gia vào cấu trúc thành tế bào và màng tế bào, quá trình phân chia tế bào và kéo dài bộ rễ. 
Magie nằm ở vị trí trung tâm của phân tử chlorophyll và có vai trò quan trọng đối với quá trình quang hợp cũng như hoạt hóa nhiều enzyme. 
Lưu huỳnh là thành phần của các amino acid chứa S như cysteine và methionine, đồng thời tham gia vào quá trình hình thành protein và chlorophyll.

Bên cạnh đó, phân đa dinh dưỡng có thể được bổ sung các nguyên tố vi lượng như B, Zn, Fe, Mn, Cu và Mo nhằm hoàn thiện thành phần dinh dưỡng của sản phẩm. 
Mặc dù cây cần các nguyên tố này với lượng rất nhỏ so với N, P và K, chúng vẫn tham gia vào nhiều quá trình sinh hóa thiết yếu và sự thiếu hụt chúng có thể trở thành yếu tố hạn chế sinh trưởng của cây.
Bo có liên quan đến tính toàn vẹn của màng, phát triển thành tế bào và quá trình hình thành ống phấn, qua đó ảnh hưởng đến khả năng hình thành hạt và quả. 
Kẽm tham gia trực tiếp hoặc gián tiếp vào nhiều hệ enzyme, quá trình tổng hợp protein, hình thành hạt và điều hòa sinh trưởng. 
Sắt có vai trò trong tổng hợp chlorophyll, hô hấp tế bào và chuyển hóa nitơ, trong khi Manganese tham gia vào hoạt hóa enzyme và quá trình quang hợp.

Các nguyên tố vi lượng có thể được đưa vào phân bón dưới nhiều dạng hóa học khác nhau, chẳng hạn \(\mathrm{ZnSO_4}\), borax, \(\mathrm{CuSO_4}\), \(\mathrm{FeSO_4}\) hoặc các dạng chelate như Zn-EDTA và Fe-EDTA.
Dạng hóa học của nguyên tố có ảnh hưởng đến tính chất sản phẩm cũng như khả năng hòa tan và sử dụng trong từng điều kiện cụ thể.

Về mặt sản xuất, phân đa dinh dưỡng có thể được tạo ra thông qua phản ứng hóa học để hình thành phân complex hoặc bằng cách phối trộn các nguyên liệu phân bón khác nhau để đạt tỷ lệ dinh dưỡng mong muốn. 
Đối với hoạt động kiểm soát chất lượng, nhóm phân đa dinh dưỡng có mức độ phức tạp cao hơn phân đơn do phải kiểm soát đồng thời nhiều chỉ tiêu. 
Hàm lượng N tổng số, \(\mathrm{P_2O_5}\), \(\mathrm{K_2O}\) và các nguyên tố trung hoặc vi lượng được công bố là cơ sở quan trọng để đánh giá chất lượng thành phẩm.
Việc phối trộn không đồng đều hoặc nguyên liệu có hàm lượng dinh dưỡng khác nhau có thể dẫn đến sự khác biệt thành phần giữa các mẫu hoặc các lô thành phẩm, vì vậy lấy mẫu đại diện và phân tích các chỉ tiêu dinh dưỡng có ý nghĩa đặc biệt trong hệ thống kiểm soát chất lượng của nhà máy.

\subsection{Phân hữu cơ}
Theo Thông tư số 07/2026/TT-BNNMT, phân bón hữu cơ là nhóm phân bón được sản xuất chủ yếu từ các chất hữu cơ tự nhiên, không bao gồm các chất hữu cơ tổng hợp, và có thể được xử lý thông qua các quá trình vật lý như làm khô, nghiền, sàng, phối trộn, làm ẩm hoặc các quá trình sinh học như ủ, lên men và chiết \cite{circular-07-2026-fertilizer}.
Nguồn nguyên liệu dùng để sản xuất phân hữu cơ có thể bao gồm phân động vật, phụ phẩm cây trồng, vật liệu thực vật, compost và các nguồn hữu cơ có khả năng cung cấp chất hữu cơ và dinh dưỡng cho đất.

Khác với phần lớn phân vô cơ, vai trò của phân hữu cơ không chỉ nằm ở việc cung cấp trực tiếp các nguyên tố dinh dưỡng mà còn ở việc bổ sung chất hữu cơ cho đất. 
Việc tăng hàm lượng chất hữu cơ có thể góp phần cải thiện cấu trúc đất, khả năng giữ nước, khả năng giữ và cung cấp dinh dưỡng cũng như hoạt động của hệ sinh vật đất.

Các chất dinh dưỡng trong vật liệu hữu cơ thường phải trải qua quá trình phân giải và khoáng hóa trước khi chuyển thành dạng mà cây có thể hấp thu. 
Vì vậy, tốc độ cung cấp dinh dưỡng của phân hữu cơ thường phụ thuộc vào đặc tính nguyên liệu, mức độ phân hủy, nhiệt độ, độ ẩm và điều kiện của đất. 
Điều này tạo ra sự khác biệt cơ bản so với nhiều loại phân vô cơ có hàm lượng dinh dưỡng cao và các chất dinh dưỡng dễ hòa tan hơn.

Một hướng phát triển quan trọng trong thực tế sản xuất là kết hợp nguồn chất hữu cơ với dinh dưỡng khoáng để tạo ra phân hữu cơ - khoáng.
Về nguyên tắc, sự kết hợp này cho phép sản phẩm vừa bổ sung chất hữu cơ cho đất vừa cung cấp các thành phần khoáng như N, P, K hoặc các nguyên tố dinh dưỡng khác theo công thức sản phẩm. 
So với phân hữu cơ thông thường, việc bổ sung nguồn dinh dưỡng khoáng vừa có thể giúp bổ sung các nguyên tố dinh dưỡng cần thiết cho cây trồng, trong khi thành phần hữu cơ vẫn đóng vai trò bổ sung vật chất hữu cơ cho đất nhằm cải thiện chất lượng đất.\cite{fao-organic-fertilizers,qcvn106-fertilizer-quality}

Một đặc điểm cần lưu ý là thành phần của phân hữu cơ có thể biến động đáng kể tùy thuộc vào nguồn nguyên liệu và phương pháp xử lý. 
Vì vậy, kiểm định nguyên liệu đầu vào và thành phẩm cũng đặc biệt quan trọng đối với sản xuất phân hữu cơ. 
FAO cho biết chất lượng phân hữu cơ có thể thay đổi rất rộng theo nguyên liệu sử dụng và đề xuất đánh giá các tiêu chí như hàm lượng chất hữu cơ, carbon hữu cơ, tổng N-\(\mathrm{P_2O_5}\)-\(\mathrm{K_2O}\) và một số kim loại nặng.\cite{fao-organic-fertilizers,qcvn106-fertilizer-quality}

\subsection{Mối liên hệ giữa nhóm phân bón và kiểm soát chất lượng}
Các nhóm phân bón khác nhau về nguồn nguyên liệu, thành phần và đặc tính dinh dưỡng nên chỉ tiêu kiểm nghiệm cũng không hoàn toàn giống nhau. Đối với phân đơn, việc kiểm định thường tập trung vào N, \(\mathrm{P_2O_5}\) hoặc \(\mathrm{K_2O}\); đối với phân đa dinh dưỡng, cần kiểm soát đồng thời nhiều chỉ tiêu đa lượng, trung lượng và vi lượng; đối với phân hữu cơ, cần quan tâm thêm đến chất hữu cơ, carbon hữu cơ và các chỉ tiêu an toàn liên quan đến nguyên liệu \cite{fao-fertilizer-specifications,qcvn106-fertilizer-quality}.

Chất lượng phân bón không chỉ phụ thuộc vào việc có mặt các nguyên tố dinh dưỡng mà còn phụ thuộc vào hàm lượng thực tế, tỷ lệ giữa các thành phần và sự ổn định giữa các lô sản xuất. Vì vậy, hoạt động lấy mẫu và phân tích nguyên liệu đầu vào, bán thành phẩm và thành phẩm của phòng thí nghiệm là mắt xích quan trọng giữa công thức sản phẩm và sản phẩm thực tế được tạo ra trên dây chuyền \cite{company-internal-documents,tcvn9486-sampling}.
